\documentclass[11pt]{article}

\usepackage[a4paper,margin=1in]{geometry}
\usepackage{amsmath,amssymb,amsfonts}
\usepackage{mathtools}
\usepackage{enumitem}
\usepackage{booktabs}
\usepackage{xcolor}
\usepackage{hyperref}

\hypersetup{
  colorlinks=true,
  linkcolor=blue!55!black,
  citecolor=blue!55!black,
  urlcolor=blue!55!black
}

\newcommand{\R}{\mathbb{R}}
\newcommand{\ud}{\boldsymbol{u}}
\newcommand{\du}{\boldsymbol{\delta u}}
\newcommand{\vv}{\boldsymbol{v}}
\newcommand{\ww}{\boldsymbol{w}}
\newcommand{\rr}{\boldsymbol{r}}
\newcommand{\HH}{\boldsymbol{H}}
\newcommand{\KK}{\boldsymbol{K}}
\newcommand{\PP}{\boldsymbol{P}}
\newcommand{\II}{\boldsymbol{I}}
\newcommand{\RR}{\boldsymbol{R}}
\newcommand{\ZZ}{\boldsymbol{Z}}
\newcommand{\MM}{\boldsymbol{M}}
\newcommand{\nn}{\boldsymbol{n}}
\newcommand{\xx}{\boldsymbol{x}}
\newcommand{\qq}{\boldsymbol{q}}
\newcommand{\dd}{\boldsymbol{d}}
\newcommand{\diag}{\operatorname{diag}}
\newcommand{\supp}{\operatorname{supp}}
\newcommand{\jump}[1]{[\![#1]\!]}
\newcommand{\negpart}[1]{\left[#1\right]_-}

\title{Filtered Nonlinear Multigrid for Nitsche Contact}
\author{Implementation note for \texttt{python/spike/mmg\_nitsche.py}}
\date{\today}

\begin{document}
\maketitle

\begin{abstract}
This note describes the nonlinear multigrid algorithm used in
\texttt{mmg\_nitsche.py} for a two-dimensional Nitsche contact problem.
The method follows the algorithmic skeleton of
\texttt{sfem\_ShiftedPenaltyMultigrid.hpp}: an outer nonlinear iteration
contains a bounded number of inner nonlinear multigrid cycles; each nonlinear
cycle performs fine-level nonlinear smoothing, residual restriction, a recursive
coarse correction, prolongation, and post-smoothing.  The current coarse space
is the jump-free active space: every Cartesian displacement component at nodes
touched by active contact rows is removed from coarse corrections.  This is the
topological, unbiased coarse active-space construction.  It is more restrictive
than a normal-jump filter because it also removes tangential components, but it
does not require normals or contact interpolation weights during the coarse
solve; only the active contact adjacency is required.
\end{abstract}

\section{Problem and notation}

Let level $\ell=0$ denote the finest mesh and let increasing $\ell$ denote
coarser levels.  The discrete displacement on level $\ell$ is
\[
    \ud_\ell \in \R^{d n_\ell},
    \qquad d=2,
\]
with prescribed Dirichlet values imposed strongly.  The finest-level nonlinear
equilibrium equation is written
\[
    \rr_0(\ud_0)
    =
    \rr_{e,0}(\ud_0)
    +
    \rr_{c,0}(\ud_0)
    =
    \boldsymbol{0},
\]
where $\rr_e$ is the elastic residual and $\rr_c$ is the Nitsche contact
residual.  The implementation uses a linear elastic tangent on the current
configuration and a semismooth active-set linearization of the contact term.
For one frozen contact geometry, the tangent is
\[
    \HH_0(\ud_0)
    =
    \HH_{e,0}
    +
    \HH_{c,0}(\ud_0).
\]
The phrase ``frozen contact geometry'' means that the contact pairing, the
surface quadrature, the closest-point interpolation, and the normal vectors are
held fixed while one residual/tangent pair is assembled.  They are recomputed
whenever the nonlinear smoother or nonlinear cycle relinearizes the contact
problem.

\section{Nitsche contact discretization}

For a potential contact row $c$, let $a(c)$ be the side-1 point and let
$\PP^c$ be the scalar interpolation row for the side-2 point.  If
$\nn^c$ is the contact normal and $\ud_1^c,\ud_2^c$ are the two traces of the
displacement, the scalar normal gap used by the code is
\[
    g^c(\ud)
    =
    \nn^{c\,T}
    \left(
        \xx_2^c + \ud_2^c
        -
        \xx_1^c - \ud_1^c
    \right).
\]
On an edge, the implementation stores a quadrature-weighted edge mean
$\bar g^c$.  With edge length $h_c$, shear modulus $\mu$, and user parameter
$\gamma_0$, the local penalty scale is
\[
    \gamma_c = \gamma_0 \frac{\mu}{h_c}.
\]
The normal elastic traction is denoted by
\[
    \sigma_n^c(\ud)
    =
    \nn^{c\,T}\boldsymbol{\sigma}(\ud)\nn^c.
\]
For the symmetric or unbiased Nitsche form, the active-set indicator is based
on the augmented normal quantity
\[
    P_n^c(\ud)
    =
    \theta_c\,\sigma_n^c(\ud)
    +
    \gamma_c\,\bar g^c(\ud),
    \qquad
    \alpha_c(\ud)
    =
    \begin{cases}
    1, & P_n^c(\ud)<0,\\
    0, & P_n^c(\ud)\ge 0.
    \end{cases}
\]
Only active rows contribute to the contact residual and tangent.  Inactive
rows are skipped.  Therefore, outside the active contact region the boundary is
natural Neumann in the discrete weak sense: no Nitsche contact force is
assembled there.  A raw plotted Cauchy value $-\sigma_n$ may still be nonzero
outside contact because it is an element stress diagnostic, not the active
contact pressure.  The active pressure diagnostic is instead
\[
    p_c(\ud)
    =
    \alpha_c(\ud)\,
    \max\!\left(-\frac{P_n^c(\ud)}{\theta_c},0\right).
\]

For a frozen active set and frozen geometry, the contact energy contribution
has the local form
\[
    \Pi_c(\ud)
    =
    \sum_{c}
    \frac{m_c}{2\,\theta_c\,\gamma_c}
    \negpart{P_n^c(\ud)}^2,
\]
where $m_c$ is the lumped contact quadrature weight.  On active rows the first
variation is
\[
    \delta \Pi_c
    =
    \sum_{c:\alpha_c=1}
    \frac{m_c}{\theta_c\gamma_c}
    P_n^c(\ud)\,
    \delta P_n^c(\ud),
\]
and the tangent contribution is
\[
    \HH_{c}
    =
    \sum_{c:\alpha_c=1}
    \frac{m_c}{\theta_c\gamma_c}
    \boldsymbol{p}_c\boldsymbol{p}_c^T,
    \qquad
    \boldsymbol{p}_c
    =
    \nabla_{\ud}P_n^c.
\]
This is exactly the structure assembled in the Python code: the derivative of
$P_n^c$ contains the gap derivative and, when enabled, the derivative of
$\sigma_n^c$.

\section{Multigrid hierarchy}

The level hierarchy consists of standard nested displacement spaces.  Let
\[
    \II_{\ell+1}^{\ell}
    :
    \R^{d n_{\ell+1}}
    \to
    \R^{d n_{\ell}}
\]
be the geometric prolongation from level $\ell+1$ to level $\ell$, and let
\[
    \RR_{\ell}^{\ell+1}
    =
    \left(\II_{\ell+1}^{\ell}\right)^T
\]
be the corresponding variational restriction.  The current Python prototype
uses inherited Galerkin coarse operators for the tangent available at the
current fine iterate:
\[
    \HH_{\ell+1}
    =
    \left(\II_{\ell+1}^{\ell}\right)^T
    \HH_{\ell}
    \II_{\ell+1}^{\ell}.
\]
This is a convenient prototype choice.  The important point for the current
experiment is not rediscretization versus inherited Galerkin elasticity, but
the filtered coarse correction space.  Once the fine tangent is fixed, the
coarse hierarchy is used as a linear correction hierarchy inside one nonlinear
cycle.

\section{Jump-free active coarse space}

\subsection{Active adjacency mask}

The implemented coarse space is the jump-free active space.  Let
$\mathcal{C}_\ell$ be the potential contact rows represented on level $\ell$,
and let $\alpha_c\in\{0,1\}$ be the frozen active indicator.  Define
$\mathcal{A}_{1,\ell}$ as the set of side-1 displacement nodes touched by
active rows, and $\mathcal{A}_{2,\ell}$ as the set of side-2 displacement nodes
touched by the interpolation stencils of active rows:
\[
    \mathcal{A}_{1,\ell}
    =
    \{a(c): c\in\mathcal{C}_\ell,\ \alpha_c=1\},
\]
\[
    \mathcal{A}_{2,\ell}
    =
    \{j: j\in\supp \PP_\ell^c,\ c\in\mathcal{C}_\ell,\ \alpha_c=1\}.
\]
For an unbiased directed contact description the side labels are local to a
row, not global labels of the physical surface.  The active displacement nodes
are therefore the union
\[
    \mathcal{A}_\ell
    =
    \mathcal{A}_{1,\ell}
    \cup
    \mathcal{A}_{2,\ell}.
\]
The inactive nodal indicator is
\[
    d_\ell^a
    =
    \begin{cases}
    0, & a\in \mathcal{A}_\ell,\\
    1, & a\notin \mathcal{A}_\ell,
    \end{cases}
\]
and the displacement filter is
\[
    \ZZ_\ell
    =
    \diag(\dd_\ell)\otimes \II_d.
\]
Applying $\ZZ_\ell$ sets both Cartesian components of active contact-adjacent
nodes to zero.  In the current implementation this is the intended starting
point: both $x$ and $y$ components are filtered.  This differs from a
normal-tangential filter, which would remove only the active normal jump and
would preserve tangential motion.

\subsection{Why the contact Hessian vanishes on the filtered space}

For an active row $c$, every displacement degree of freedom participating in
the row is masked by $\ZZ_\ell$.  Hence, for any vector $\vv_\ell$,
\[
    \left(\ZZ_\ell\vv_\ell\right)_{1}^{c}
    =
    \boldsymbol{0},
    \qquad
    \left(\ZZ_\ell\vv_\ell\right)_{2}^{j}
    =
    \boldsymbol{0}
    \quad
    \forall j\in\supp \PP_\ell^c.
\]
Therefore the active normal jump induced by the filtered vector is zero:
\[
    \nn^{c\,T}
    \left[
        \left(\ZZ_\ell\vv_\ell\right)_1^c
        -
        \sum_{j\in\supp \PP_\ell^c}
        P_{\ell,cj}
        \left(\ZZ_\ell\vv_\ell\right)_2^j
    \right]
    =
    0.
\]
It follows that the active contact tangent has no action on the filtered
coarse correction space:
\[
    \ZZ_\ell^T
    \HH_{c,\ell}
    \ZZ_\ell
    =
    \boldsymbol{0},
\]
provided the active set and contact geometry are frozen during the correction.
Consequently, the coarse correction can be driven by the elastic part of the
operator restricted to the jump-free space.  In the current prototype, because
the whole fine tangent is inherited by Galerkin projection, the same effect is
enforced algebraically by zeroing filtered degrees of freedom in the coarse
right-hand side, coarse solve, and prolongated correction.

\subsection{Coarse mask transfer}

The fine active mask is built from the current fine active contact rows.  On
coarser levels, the current implementation restricts the active nodal mask
through the variational support of the geometric transfer and thresholds the
result:
\[
    s_{\ell+1}
    =
    \chi
    \left(
        \left(\II_{\ell+1}^{\ell}\right)^T s_{\ell}
    \right),
    \qquad
    d_{\ell+1}=1-s_{\ell+1}.
\]
Here $\chi(z)=1$ for positive entries and $\chi(z)=0$ otherwise.  Dirichlet
degrees of freedom are always included in the constrained mask.  The complete
mask on level $\ell$ is therefore
\[
    \mathcal{M}_\ell
    =
    \mathcal{D}_\ell
    \cup
    \{(a,i): a\in\mathcal{A}_\ell,\ i=1,\ldots,d\}.
\]

\section{Recursive linear correction cycle}

Inside one nonlinear cycle, after the fine tangent and masks have been built,
the coarse solve is a linear multigrid correction.  Given a level $\ell$,
right-hand side $\boldsymbol{b}_\ell$, operator $\HH_\ell$, and mask
$\mathcal{M}_\ell$, the masked right-hand side is
\[
    \widehat{\boldsymbol{b}}_\ell
    =
    \ZZ_\ell^T\boldsymbol{b}_\ell,
\]
with all masked entries set to zero.  The correction is constrained by
\[
    \du_\ell\big|_{\mathcal{M}_\ell}=0.
\]
On the coarsest level the implementation solves the masked system directly on
the unmasked degrees of freedom.  On intermediate levels the recursive cycle is
\[
    \du_\ell^{(0)} = \boldsymbol{0},
\]
\[
    \du_\ell^{(1)}
    =
    S_\ell^{\nu_1}
    \left(
        \widehat{\boldsymbol{b}}_\ell,
        \du_\ell^{(0)}
    \right),
\]
\[
    \boldsymbol{\rho}_\ell
    =
    \widehat{\boldsymbol{b}}_\ell
    -
    \HH_\ell\du_\ell^{(1)},
    \qquad
    \boldsymbol{\rho}_\ell\big|_{\mathcal{M}_\ell}=0,
\]
\[
    \boldsymbol{b}_{\ell+1}
    =
    \left(\II_{\ell+1}^{\ell}\right)^T
    \boldsymbol{\rho}_\ell,
\]
\[
    \du_{\ell+1}
    =
    \operatorname{cycle}_{\ell+1}
    \left(
        \boldsymbol{b}_{\ell+1}
    \right),
\]
\[
    \du_\ell^{(2)}
    =
    \du_\ell^{(1)}
    +
    \II_{\ell+1}^{\ell}\du_{\ell+1},
    \qquad
    \du_\ell^{(2)}\big|_{\mathcal{M}_\ell}=0,
\]
\[
    \du_\ell
    =
    S_\ell^{\nu_2}
    \left(
        \widehat{\boldsymbol{b}}_\ell,
        \du_\ell^{(2)}
    \right).
\]
The parameter \texttt{cycle\_type} repeats the coarse-correction part once for
a V-cycle and twice for a W-cycle, matching the structure of
\texttt{ShiftedPenaltyMultigrid::cycle}.  The default Python smoother is
Jacobi; Gauss--Seidel, symmetric Gauss--Seidel, and block Jacobi are also
available for experiments.

\section{Fine nonlinear smoothing}

Fine smoothing is nonlinear because each smoothing step reassembles the
current contact residual and tangent.  Let $S_0$ denote one linear smoother
application to the current tangent.  One nonlinear smoothing step is
\[
    \rr_0^k = \rr_0(\ud_0^k),
    \qquad
    \HH_0^k = \HH_0(\ud_0^k),
\]
\[
    \HH_0^k \boldsymbol{s}^k
    \approx
    -\rr_0^k,
    \qquad
    \ud_0^{k+1}
    =
    \ud_0^k + \boldsymbol{s}^k,
\]
where the approximate solve is a fixed number of masked smoother sweeps.  Only
Dirichlet degrees of freedom are fixed during fine smoothing.  Active contact
nodes are allowed to move on the fine level; the jump-free active mask is used
for coarse corrections, not for suppressing fine-level nonlinear relaxation.

\section{Nonlinear multigrid cycle}

One nonlinear multigrid cycle at the finest level has the following structure.
The active set, masks, residual, and tangent are refreshed before the cycle and
again after each nonlinear smoothing phase.

\begin{enumerate}[label=\textbf{Step \arabic*.},leftmargin=2.2cm]
\item \textbf{Refresh contact state.}
Compute $\rr_0(\ud_0)$, $\HH_0(\ud_0)$, the active contact set
$\mathcal{A}_0$, and the hierarchy masks $\mathcal{M}_\ell$.

\item \textbf{Pre-smooth.}
Apply $\nu_{\mathrm{nl}}$ nonlinear smoothing steps on the fine level, each
using $\nu_1$ linear smoother sweeps.

\item \textbf{Restrict residual.}
Recompute the fine residual and tangent.  Form
\[
    \boldsymbol{b}_1
    =
    -\left(\II_1^0\right)^T\rr_0(\ud_0),
\]
with fine Dirichlet rows zeroed before restriction.

\item \textbf{Coarse correction.}
Build the inherited Galerkin hierarchy from the current fine tangent and solve
approximately
\[
    \HH_1 \du_1
    =
    \boldsymbol{b}_1
\]
with the recursive masked linear cycle.

\item \textbf{Prolong and filter.}
Apply the correction
\[
    \du_0
    =
    \II_1^0 \du_1,
    \qquad
    \du_0\big|_{\mathcal{M}_0}=0,
\]
and update
\[
    \ud_0 \leftarrow \ud_0 + \du_0.
\]

\item \textbf{Post-smooth.}
Apply $\nu_{\mathrm{nl}}$ nonlinear smoothing steps on the fine level, each
using $\nu_2$ linear smoother sweeps.
\end{enumerate}

If \texttt{--skip-coarse} is enabled, Step 4 and Step 5 are omitted.  This is
useful for isolating whether the smoother alone reduces the residual.

\section{Outer and inner iteration}

The top-level nonlinear solve follows the same outer/inner structure as the
shifted-penalty multigrid implementation.  Let $m_{\max}$ be the maximum
number of outer iterations and $q_{\max}$ the maximum number of inner
nonlinear cycles per outer iteration.  Starting from the Dirichlet displacement
field, the solver computes
\[
    r_0 = \|\rr_0(\ud_0^0)\|_2
\]
and then performs
\[
    \text{for } k=0,\ldots,m_{\max}-1:
    \qquad
    \text{for } q=0,\ldots,q_{\max}-1:
    \qquad
    \ud_0 \leftarrow \operatorname{NMG}(\ud_0).
\]
After each inner nonlinear cycle, the residual norm
\[
    r_{k,q} = \|\rr_0(\ud_0)\|_2
\]
and relative residual $r_{k,q}/r_0$ are checked.  The default residual target
is
\[
    r_{k,q}<10^{-10}
    \quad\text{or}\quad
    \frac{r_{k,q}}{r_0}<10^{-10}.
\]
The implementation also monitors the active penetration norm
\[
    \|g_-\|_2
    =
    \left\|
    \min(g(\ud_0),0)
    \right\|_2,
\]
and reports active rows, contact quadrature points, filtered degrees of
freedom, and contact tangent sparsity.  A stagnation guard terminates an inner
loop when the residual reduction within an outer iteration becomes too small.

\section{Algorithm summary}

\begin{center}
\begin{minipage}{0.94\linewidth}
\hrule
\vspace{0.5em}
\textbf{Algorithm: filtered nonlinear multigrid for Nitsche contact}
\begin{enumerate}[leftmargin=1.7em]
\item Build the geometric hierarchy and prolongation operators.
\item Initialize the fine displacement with the prescribed Dirichlet values.
\item Compute the initial residual norm $r_0$.
\item For each outer iteration:
  \begin{enumerate}[leftmargin=1.7em]
  \item For each inner nonlinear cycle:
    \begin{enumerate}[leftmargin=1.7em]
    \item Assemble the fine Nitsche residual and tangent with frozen geometry.
    \item Build the active contact adjacency mask and restrict it through the
    hierarchy.
    \item Apply fine nonlinear pre-smoothing.
    \item Reassemble and restrict the current residual.
    \item Compute a masked recursive coarse correction.
    \item Prolong the correction, zero the fine active-mask entries, and update
    the fine displacement.
    \item Apply fine nonlinear post-smoothing.
    \item Check residual convergence or stagnation.
    \end{enumerate}
  \item Report residual, penetration, active set size, and filtering
  diagnostics.
  \end{enumerate}
\item Stop when the residual and penetration criteria are satisfied or when the
maximum iteration count is reached.
\end{enumerate}
\vspace{0.5em}
\hrule
\end{minipage}
\end{center}

\section{Relation to the normal-jump filter}

The jump-free filter $\ZZ_\ell$ is intentionally more restrictive than the
active normal-jump filter.  A normal-jump filter would work in
normal--tangential coordinates and remove only
\[
    \nn^{c\,T}\jump{\du_\ell}^c
\]
for active contact rows, while leaving tangential corrections available.  Such
a filter requires the coarse contact interpolation, coarse active rows, and
coarse normals.  Its advantage is that it preserves more interface modes.  Its
cost is that the multigrid transfer must correctly evaluate and remove the
normal jump on every level.

The current prototype deliberately starts with the simpler unbiased
jump-free construction:
\[
    \ZZ_\ell=\diag(\dd_\ell)\otimes\II_d.
\]
This filter does not distinguish normal and tangential components; it removes
both $x$ and $y$ displacement components at active-adjacent nodes.  It is
therefore a robust first test of the multigrid skeleton and the smoother.  If
the method later needs a less restrictive coarse space, the natural extension
is to replace $\ZZ_\ell$ by a normal--tangential active-jump filter constructed
from level-wise normals and contact interpolation rows.

\section{Implementation status and diagnostics}

The current Python implementation has the following important properties.

\begin{itemize}[leftmargin=1.7em]
\item The nonlinear-cycle structure matches the shifted-penalty multigrid
skeleton: pre-smooth, restrict residual, recursive coarse solve, prolong,
post-smooth.
\item The coarse masks are refreshed from the current fine active set at each
nonlinear cycle.  The Galerkin operator itself needs no independent contact
refresh beyond the current residual/tangent linearization and the active-space
mask.
\item Fine smoothing moves active contact nodes; coarse correction does not.
\item The default smoother is scalar Jacobi.  It was selected because, in the
current test case, it had better late residual reduction than forward
Gauss--Seidel.
\item The active contact pressure is the masked Nitsche quantity
$\alpha_c\max(-P_n^c/\theta_c,0)$, not the raw Cauchy stress diagnostic
$-\sigma_n^c$.
\end{itemize}

\section{Expected failure modes}

The residual can stagnate or decay slowly for several reasons.

\begin{enumerate}[leftmargin=1.7em]
\item \textbf{Overly restrictive coarse space.}
The jump-free filter removes tangential as well as normal components at active
nodes.  If the active union mask becomes large, the coarse correction may lose
too much interface content.

\item \textbf{Insufficient smoothing.}
The contact tangent is localized and stiff.  Too few smoothing sweeps can leave
high-frequency error that the coarse grid cannot represent.

\item \textbf{Changing active set.}
The coarse correction assumes frozen contact geometry and a frozen active set.
If the active set changes significantly between nonlinear cycles, the inherited
coarse correction may become a poor model of the next residual.

\item \textbf{Diagnostic confusion.}
Raw $-\sigma_n$ may be nonzero outside contact even when the inactive boundary
is natural Neumann in the weak sense.  Solver convergence should be assessed
with the residual norm, penetration, and active Nitsche pressure.

\item \textbf{Coordinate-space mismatch in future filters.}
If the method is changed from the jump-free filter to a normal-jump filter, the
restriction and prolongation must be composed with the correct
normal--tangential transformations.  Removing a Cartesian component is not the
same as removing the active normal component.
\end{enumerate}

\section{Conclusion}

The implemented method is a filtered nonlinear multigrid solver for Nitsche
contact.  It currently uses a conservative jump-free coarse active space:
coarse corrections are not allowed to move displacement nodes touched by active
contact rows.  This makes the active contact Hessian vanish on the coarse
correction space and leaves the recursive multigrid cycle to reduce the
elastic bulk error away from the rapidly changing contact interface.  The
algorithm is deliberately structured like the existing shifted-penalty
multigrid solver, while using the contact-specific residual, active set, and
filtered coarse-space construction needed for the Nitsche formulation.

\end{document}
